\chapter{Rachunek całek cechowanych}
\providecommand{\atlasref}[2]{#2}
\label{ma1:ch:rachunek-calek-cechowanych}
\begingroup
\providecommand{\osc}{\operatorname{osc}}
\providecommand{\Var}{\operatorname{Var}}
\providecommand{\mesh}{\operatorname{mesh}}
\providecommand{\supp}{\operatorname{supp}}
\providecommand{\AC}{\mathrm{AC}}
\providecommand{\BV}{\mathrm{BV}}
\providecommand{\HK}{\mathrm{KH}}
\providecommand{\ind}[1]{\chi_{#1}}
\providecommand{\Dn}{\underline D}
\providecommand{\Du}{\overline D}
\providecommand{\KS}{\mathrm{KS}}
Przyjmujemy oznaczenia i konwencje z rozdziału \ref{ma1:ch:kurzweil-henstock}.
\section{Podstawienie oraz całkowanie przez części}
\begin{twierdzenie}[Mnożniki o ograniczonym wahaniu]\label{ma1:kh:thm:bv-multiplier}
Niech $f\in\HK[a,b]$, $F(x)=\int_a^x f$, a $g:[a,b]\to\R$ ma ograniczone wahanie. Wtedy $fg\in\HK[a,b]$ oraz
\[
\int_a^b f(x)g(x)\,\mathrm dx=F(b)g(b)-F(a)g(a)-\int_a^b F\,\mathrm dg,
\]
gdzie ostatnia całka jest zwykłą całką Riemanna--Stieltjesa.
\end{twierdzenie}
\begin{proof}
Dla $a=b$ obie strony są zerowe. Niech $a<b$.
Z twierdzenia \ref{ma1:kh:thm:kh-primitive-derivative}
funkcja $F$ jest ciągła, a $g$ jest regulowana na mocy
lematu \ref{it:lem:ks-bv-regulated}. Ponieważ $g\in\BV[a,b]$,
twierdzenie \ref{tw:ks-bv-reg} zapewnia istnienie całki
$(\KS)\int_a^b g\,\mathrm dF$.
Dla integratora $\alpha(x)=x$ całka KS jest dokładnie
całką KH. Twierdzenie o podstawieniu całki nieoznaczonej
\ref{tw:ks-chain} daje więc $fg\in\HK[a,b]$ oraz
\[
 (\HK)\int_a^b fg=(\KS)\int_a^b g\,\mathrm dF.
\]
Stosując wzór całkowania przez części
\ref{tw:ks-ibp}, w którym wszystkie poprawki skokowe
znikają dzięki ciągłości $F$, otrzymujemy
\[
 (\KS)\int_a^b g\,\mathrm dF
 =F(b)g(b)-F(a)g(a)-(\KS)\int_a^b F\,\mathrm dg.
\]
Z twierdzenia \ref{it:thm:rs-tools}(a) istnieje
zwykła całka Riemanna--Stieltjesa z $F$ względem $g$.
Twierdzenie \ref{tw:ks-cauchy} daje jej zgodność
z ostatnią całką KS, co kończy dowód.
\end{proof}
\begin{uwaga}[Ograniczenia całki KH]
Klasa $\HK[a,b]$ nie jest kratą: z całkowalności $f$ nie wynika całkowalność $|f|$ (przykład \ref{ma1:kh:ex:conditional-derivative}). Nie jest też algebrą, ponieważ $x^{-1/2}$ na $(0,1]$ po nadaniu wartości zero w zerze jest całkowalna
w sensie Lebesgue'a, a więc także KH
(twierdzenie \ref{it:thm:mcshane-lebesgue} i inkluzja
McShane'a w KH). Jej kwadrat nie jest KH-całkowalny,
gdyż jest nieujemny i nie jest całkowalny w sensie
Lebesgue'a (twierdzenie \ref{ma1:kh:thm:kh-absolute}). Dla $a<b$ półnorma Alexiewicza $\|f\|_A=\sup_{x\in[a,b]}\left|\int_a^x f\right|$ jest normą po utożsamieniu funkcji równych prawie wszędzie, lecz otrzymana przestrzeń nie jest zupełna.
\end{uwaga}
\begin{twierdzenie}[Podstawienie przez dyfeomorfizm klasy $C^1$]\label{ma1:kh:tw:hk-podstawienie}\label{tw:hk-podstawienie}
Niech $a<b$, $u<v$, $f\in\HK[a,b]$ i niech $\phi:[u,v]\to[a,b]$ będzie rosnącym dyfeomorfizmem klasy $C^1$, to znaczy $\phi\in C^1[u,v]$, $\phi(u)=a$, $\phi(v)=b$ i $\phi'(t)>0$ dla $t\in[u,v]$. Wówczas $t\mapsto f(\phi(t))\phi'(t)$ jest KH-całkowalna na $[u,v]$ oraz
\[
\int_u^v f(\phi(t))\phi'(t)\,\mathrm dt=\int_a^b f(x)\,\mathrm dx.
\]
Dla malejącego dyfeomorfizmu z $\phi(u)=b$, $\phi(v)=a$ i $\phi'<0$ całka po prawej ma odwrócone granice.
\end{twierdzenie}
\begin{proof}
Niech $F(x)=\int_a^x f$. Z twierdzenia \ref{ma1:kh:thm:kh-primitive-derivative}
funkcja $F$ jest ciągła i $F'=f$ prawie wszędzie. Najpierw uzasadnimy, że $F\circ\phi$ należy do $\AC G_*[u,v]$. Z charakteryzacji pierwotnych KH (twierdzenie \ref{ma1:kh:thm:kh-acgstar-characterization}) mamy $[a,b]=\bigcup_n E_n$, gdzie zbiory $E_n$ są domknięte i $F\in\AC_*(E_n)$. Zbiory $\phi^{-1}(E_n)$ są domknięte i pokrywają $[u,v]$. Ustalmy $n$ i $\varepsilon>0$. Z warunku $\AC_*(E_n)$
wybieramy $\eta>0$ dla błędu $\varepsilon$.
Niech $M=\max_{[u,v]}\phi'>0$. Z twierdzenia Lagrange'a
$\phi$ jest $M$-lipschitzowska. Dla przedziałów $J_j$
o rozłącznych wnętrzach, końcach w $\phi^{-1}(E_n)$
i $\sum_j|J_j|<\eta/M$, ich obrazy mają rozłączne
wnętrza, końce w $E_n$ oraz $\sum_j|\phi(J_j)|<\eta$.
Równość oscylacji daje
\[
 \sum_j\osc_{J_j}(F\circ\phi)=\sum_j\osc_{\phi(J_j)}F
 <\varepsilon.
\]
Zatem $F\circ\phi\in\AC_*(\phi^{-1}(E_n))$,
a ciągłość złożenia i pokrycie domknięte dają
$F\circ\phi\in\AC G_*[u,v]$.

Z twierdzenia Lagrange'a, dla $m=\min_{[u,v]}\phi'>0$ oraz zdefiniowanego wyżej $M$, mamy $m|s-t|\le|\phi(s)-\phi(t)|\le M|s-t|$. Dyfeomorfizm i jego odwrotność są więc lipschitzowskie. Dla zbioru zerowego $N\subset[a,b]$ i $\varepsilon>0$
wybieramy przeliczalne pokrycie przedziałami $I_j$
o $\sum_j|I_j|<m\varepsilon$. Przeciwobrazy
$\phi^{-1}(I_j\cap[a,b])$ są przedziałami
(pomijamy puste); z lipschitzowskości odwrotności
suma ich długości jest mniejsza od $\varepsilon$. Przeciwobraz jest zatem zbiorem miary zero. W punktach, w których $F'(\phi(t))=f(\phi(t))$, zwykła reguła łańcuchowa daje $(F\circ\phi)'(t)=f(\phi(t))\phi'(t)$. Twierdzenie \ref{ma1:kh:thm:acgstar-derivative-integrable} daje całkowalność KH funkcji $t\mapsto f(\phi(t))\phi'(t)$, gdyż ma skończone wartości i zgadza się z pochodną pierwotnej prawie wszędzie, oraz wzór Newtona--Leibniza; ostatni przyrost wynosi $F(b)-F(a)$. Dla malejącego dyfeomorfizmu argument jest taki sam: obrazy komórek zapisujemy z końcami w kolejności rosnącej, ich długości kontroluje $\max|\phi'|$, a odwrotność jest lipschitzowska z oszacowania $\min|\phi'|>0$. Przyrost funkcji złożonej wynosi tym razem $F(a)-F(b)$.
\end{proof}
\begin{uwaga}[Zakres twierdzenia]
Warunek $\phi'(t)>0$ zabezpiecza przed trudnościami z przenoszeniem zbiorów wyjątkowych. Istnieją mocniejsze twierdzenia o podstawieniu dla mniej regularnych funkcji $\phi$, ale bez dodatkowych warunków na składanie pierwotnych i na zbiory wyjątkowe nie wolno automatycznie zastępować dyfeomorfizmu dowolną funkcją absolutnie ciągłą.
\end{uwaga}
\begin{wniosek}[Całkowanie przez części z mnożnikiem absolutnie ciągłym]\label{ma1:kh:wn:hk-ibp-ac}\label{wn:hk-ibp-ac}
Niech $f\in\HK[a,b]$, $F(x)=\int_a^x f$, a $g\in\AC[a,b]$. Wtedy $fg\in\HK[a,b]$ i
\[
\int_a^b f(x)g(x)\,\mathrm dx=F(b)g(b)-F(a)g(a)-\int_a^b F(x)g'(x)\,\mathrm d\lambda_1(x).
\]
Ostatnia całka jest całką Lebesgue'a, ponieważ $F$ jest ograniczona i $g'\in L^1[a,b]$.
\end{wniosek}
\begin{proof}
Dla $a=b$ teza jest bezpośrednia.
Dla $a<b$ funkcja $g$ ma ograniczone wahanie,
a dokładnie $V_a^b(g)=\int_{[a,b]}|g'|\,\mathrm d\lambda_1$
(twierdzenie \ref{ma1:thm:ac-banach-separable}). Stosujemy twierdzenie \ref{ma1:kh:thm:bv-multiplier}. Z twierdzenia \ref{ma1:kh:thm:ac-lebesgue-primitive} mamy $g(v)-g(u)=\int_u^v g'\,\mathrm d\lambda_1$. Różnica sumy Stieltjesa i całki Lebesgue'a jest zatem ograniczona przez
\[
 \sum_i\int_{x_{i-1}}^{x_i}|F(t_i)-F(x)|\,|g'(x)|\,\mathrm d\lambda_1(x)
 \le \omega_F(\mesh P)\int_a^b|g'|\,\mathrm d\lambda_1.
\]
Prawa strona dąży do zera, co daje żądaną identyfikację.
\end{proof}
\begin{przyklad}[Mnożnik i całka niewłaściwa]
Niech $f(x)=\sin(1/x)/x$ dla $x\in(0,1]$ i $f(0)=0$.
Z twierdzenia \ref{ma1:kh:tw:hk-oscylacyjne} dla $p=q=1$
mamy $f\in\HK[0,1]\setminus\mathcal L^1[0,1]$,
a podstawienie $t=1/x$ na $[\varepsilon,1]$
daje $\int_\varepsilon^1f=\int_1^{1/\varepsilon}(\sin t)/t\,\mathrm dt$.
Dla mnożnika $g(x)=x$ iloczyn $fg$ jest ograniczony
i mierzalny, więc jest całkowalny nawet w sensie Lebesgue'a.
Wzór z wniosku \ref{ma1:kh:wn:hk-ibp-ac} przy $g'=1$
daje ponadto
\[
 (\HK)\int_0^1 x f(x)\,\mathrm dx
 =F(1)-\int_{[0,1]} F(x)\,\mathrm d\lambda_1(x),
 \qquad F(x)=(\HK)\int_0^x f.
\]
W ten sposób całkę iloczynu wyrażamy za pomocą
wartości i zwykłej całki z ciągłej funkcji pierwotnej.
\end{przyklad}

\begin{stwierdzenie}[Uzupełnienie w normie Alexiewicza]\label{it:thm:alexiewicz-completion}
Niech $a<b$. Przestrzeń funkcji KH-całkowalnych modulo równość prawie wszędzie, z normą $\|[f]\|_A=\sup_{x\in[a,b]}|(\HK)\int_a^x f|$, nie jest zupełna. Jej uzupełnienie izometryczne można utożsamić z $C_0[a,b]=\{F\in C[a,b]:F(a)=0\}$ z normą supremum (tutaj $C_0[a,b]$ oznacza funkcje zerujące się w $a$, nie przestrzeń funkcji znikających w nieskończoności).
\end{stwierdzenie}
\begin{proof}
Przyporządkowanie $f\mapsto F(x)=(\HK)\int_a^x f$ przyjmuje wartości w $C_0[a,b]$ na mocy twierdzenia \ref{ma1:kh:thm:kh-primitive-derivative} o ciągłości pierwotnej. Jest liniowe dzięki liniowości całki na każdym
podprzedziale (lokalność: twierdzenie
\ref{ma1:kh:thm:kh-locality}).
Jeżeli $f=0$ prawie wszędzie, lemat
\ref{ma1:kh:lem:kh-null-modification} zastosowany na
każdym $[a,x]$ daje $F(x)=0$. Odwrotnie, jeśli $F=0$,
to $f=F'=0$ prawie wszędzie na mocy twierdzenia
\ref{ma1:kh:thm:kh-primitive-derivative}.
Jądrem odwzorowania jest więc dokładnie podprzestrzeń
funkcji zerowych prawie wszędzie. Na przestrzeni ilorazowej jest zatem różnowartościową izometrią, ponieważ $\|F\|_\infty=\|[f]\|_A$. W szczególności $\|\cdot\|_A$ jest na niej normą. Gęstość obrazu wykażemy bezpośrednio: dla $F\in C_0[a,b]$ i podziału jednostajnego na $n$ części niech $P_n$ będzie interpolantem liniowym wartości $F$ w węzłach. Mamy $P_n(a)=0$ oraz $\|P_n-F\|_\infty\le\omega_F((b-a)/n)\to0$, gdzie $\omega_F(r)=\sup\{|F(x)-F(y)|:x,y\in[a,b],\ |x-y|\le r\}$. Oszacowanie wynika z faktu, że na każdej komórce $P_n(x)$ jest kombinacją wypukłą wartości $F$ na jej końcach; zbieżność do zera wynika z jednostajnej ciągłości $F$ na przedziale zwartym. Funkcja nachyleń $p_n$, stała we wnętrzach komórek i równa zeru w węzłach, jest całkowalna w sensie Riemanna, a $\int_a^x p_n=P_n(x)$ przez sumowanie przyrostów liniowych. Zatem $P_n$ należy do obrazu, a funkcje te są gęste w $C_0[a,b]$, więc domknięciem obrazu jest całe $C_0[a,b]$. Nie każda funkcja ciągła o wartości zero w $a$ jest pierwotną KH. Niech $C$ będzie funkcją Cantora z przykładu
\atlasref{ME-17::ex:cantor-measure-17}{ME-17: miara Cantora}.
Funkcja $F(x)=C((x-a)/(b-a))$ jest ciągła, $F(a)=0$, $F(b)=1$.
Poza afinicznym przeciwobrazem zbioru Cantora, który ma miarę zero
na mocy \atlasref{ME-13::eq:linear-volume-13}{ME-13: zmiana objętości},
funkcja jest lokalnie stała. Gdyby była pierwotną KH funkcji $f$,
twierdzenie \ref{ma1:kh:thm:kh-primitive-derivative}
dawałoby $f=F'=0$ prawie wszędzie. Lemat
\ref{ma1:kh:lem:kh-null-modification} zastosowany na każdym $[a,x]$
dawałby $F(x)=0$, sprzecznie z $F(b)=1$.
Przestrzeń $C[a,b]$ z normą supremum jest przestrzenią Banacha
(por. \atlasref{FA-12::prop:ck-banach12}{FA-12: zupełność przestrzeni $C(K)$}).
Przestrzeń $C_0[a,b]$ jest jej domkniętą podprzestrzenią:
jest jądrem ciągłego funkcjonału ewaluacji $F\mapsto F(a)$,
dla którego $|F(a)|\le\|F\|_\infty$. Jest więc zupełna. Obraz jest gęsty, lecz nie jest całym $C_0[a,b]$, a więc nie jest domknięty i nie jest zupełny.
\end{proof}

\section{Twierdzenia o zbieżności}
\begin{twierdzenie}[Zbieżność z majorantą Lebesgue'a]\label{ma1:kh:tw:hk-zbieznosc-dct}\label{tw:hk-zbieznosc-dct}
Niech $f_n:[a,b]\to\R$ będą mierzalne i KH-całkowalne, $f_n\to f$ prawie wszędzie dla $f:[a,b]\to\R$, a $g\in\mathcal L^1[a,b]$ niech spełnia $|f_n|\le g$ prawie wszędzie dla każdego $n$. Wówczas $f,f_n\in\mathcal L^1[a,b]$ i
$\lim_{n\to\infty}\int_a^bf_n\,\mathrm dx=\int_a^bf\,\mathrm dx$, przy czym po obu stronach całkę można rozumieć zarówno w sensie Lebesgue'a, jak i KH.
\end{twierdzenie}
\begin{proof}
Przeliczalna suma zbiorów wyjątkowych dla nierówności i zbieżności jest zbiorem zerowym. Poza nim $f$ jest granicą funkcji mierzalnych, a więc z zupełności miary Lebesgue'a także $f$ jest mierzalna. Z oszacowania $|f_n|\le g$ wynika całkowalność w sensie Lebesgue'a każdej $f_n$; zbieżność prawie wszędzie i $|f|\le g$ dają to samo dla $f$. Stosujemy twierdzenie \atlasref{ME-08::thm:dct}{ME-08: twierdzenie o zbieżności zmajoryzowanej} oraz zgodność całek Lebesgue'a i KH (twierdzenie \ref{it:thm:mcshane-lebesgue} i inkluzja całki McShane'a w KH). Nieujemna majoranta KH-całkowalna jest automatycznie
całkowalna w sensie Lebesgue'a na mocy twierdzenia
\ref{ma1:kh:thm:kh-absolute}. Dwustronne ograniczenie przez funkcje KH-całkowalne, które same nie muszą należeć do $L^1$, oraz zbieżność monotoniczną omawia twierdzenie \ref{ma1:kh:thm:kh-order-limits}.
\end{proof}
\begin{przyklad}[Dlaczego sama zbieżność punktowa nie wystarcza]
Niech $f_n(0)=0$ i $f_n(x)=n\chi_{(0,1/n]}(x)$ na $[0,1]$. Każda $f_n$ jest całkowalna w sensie Riemanna, a więc KH-całkowalna, oraz $f_n(x)\to0$ w każdym punkcie $[0,1]$. Mimo to $\int_0^1f_n\,\mathrm dx=1$ dla każdego $n$. Nie istnieje jedna całkowalna w sensie Lebesgue'a funkcja $g$ z $|f_n|\le g$ prawie wszędzie dla wszystkich $n$: w przeciwnym razie twierdzenie o zbieżności zmajoryzowanej wymuszałoby granicę całek równą zeru.
\end{przyklad}
\begin{stwierdzenie}[Jednostajna zbieżność]\label{ma1:kh:st:hk-zbieznosc-jednostajna}\label{st:hk-zbieznosc-jednostajna}
Jeśli $f_n\in\HK[a,b]$ i $f_n\to f$ jednostajnie, to $f\in\HK[a,b]$ oraz $\int_a^bf_n\to\int_a^bf$.
\end{stwierdzenie}
\begin{proof}
Stosujemy twierdzenie \ref{it:thm:ks-uniform}
dla integratora $g(x)=x$, którego wahanie wynosi $b-a$.
Sumy KS dla tego integratora są sumami KH, zatem
całkowalność i zbieżność wartości wynikają bezpośrednio
z tego twierdzenia. Nie wymagamy ograniczoności samych
$f_n$ ani $f$.
\end{proof}


\begin{twierdzenie}[Jednostajna całkowalność względem cechowania]
\label{ma1:kh:thm:kh-equiintegrability}\label{it:thm:kh-equiintegrability}
Niech $f_n\in\HK[a,b]$ i $f_n(t)\to f(t)\in\R$ dla każdego $t\in[a,b]$. Załóżmy, że dla każdego $\varepsilon>0$ istnieje jedno cechowanie $\delta$ takie, że
\[
 |S(f_n,\mathcal P)-I_n|<\varepsilon
 \quad\text{dla wszystkich }n\ge1
 \text{ i wszystkich }\delta\text{-drobnych }\mathcal P,
 \qquad I_n=\int_a^b f_n.
\]
Wtedy $f\in\HK[a,b]$ oraz $I_n\to\int_a^b f$.
\end{twierdzenie}
\begin{proof}
Dla $g(x)=x$ mamy $S_g(f_n,\mathcal P)=S(f_n,\mathcal P)$.
Przyjęte założenie o wspólnym cechowaniu jest dokładnie
założeniem twierdzenia \ref{it:thm:ks-equiintegrability}.
Stosując je do tego integratora, otrzymujemy
$f\in\HK[a,b]$ i zbieżność całek, bez uprzedniego
zakładania całkowalności funkcji granicznej.
\end{proof}
\begin{twierdzenie}[Twierdzenia graniczne dla KH]\label{ma1:kh:thm:kh-order-limits}\label{it:thm:hk-limit-theorems}
Niech $(f_n)$ będzie ciągiem rzeczywistych funkcji
KH-całkowalnych na $[a,b]$, a $f:[a,b]\to\R$.
W przypadku granicy określonej tylko prawie wszędzie
wybieramy dowolne skończone wartości na zbiorze
wyjątkowym. Wszystkie całki bez oznaczenia miary
w poniższych trzech warunkach rozumiemy w sensie KH.
\begin{enumerate}[label=(\alph*)]
\item Jeśli $f_n\uparrow f$ prawie wszędzie i $\sup_n(\HK)\int_a^b f_n<\infty$, to $f$ jest KH-całkowalna oraz $\int_a^b f_n\to\int_a^b f$.
\item Jeśli $f_n\to f$ prawie wszędzie i istnieją $g,h\in\HK[a,b]$ takie, że $g\le f_n\le h$ prawie wszędzie dla wszystkich $n$, to $f\in\HK[a,b]$ i $\int_a^b f_n\to\int_a^b f$.
\item Jeśli $f_n\downarrow f$ prawie wszędzie i
$\inf_n(\HK)\int_a^b f_n>-\infty$, to $f\in\HK[a,b]$
i całki zbiegają do całki $f$.
\end{enumerate}
We wszystkich trzech przypadkach $f_n-f\in\mathcal L^1$
oraz dla pierwotnych $F_n(x)=(\HK)\int_a^x f_n$,
$F(x)=(\HK)\int_a^x f$ zachodzi
\[
 \|[f_n-f]\|_A
 =\|F_n-F\|_\infty
 \le V_a^b(F_n-F)
 =\int_{[a,b]}|f_n-f|\,\mathrm d\lambda_1\longrightarrow0.
\]
Nie zakładamy, że same $f_n$ lub $f$ są całkowalne
w sensie Lebesgue'a.
\end{twierdzenie}
\begin{proof}
Jeżeli $a=b$, całki oraz wahania są zerowe, a wszystkie
funkcje należą do $\mathcal L^1$ względem miary zerowej.
Dalej niech $a<b$. Dla integratora $\alpha(x)=x$ mamy
$\mu_\alpha=\lambda_1|_{[a,b]}$ na zbiorach borelowskich,
ponieważ obie miary przypisują $[a,t]$ masę $t-a$
(twierdzenie \ref{tw:ks-measure}). Ich uzupełnienia
również są równe. Funkcje KH-całkowalne są mierzalne
w sensie Lebesgue'a na mocy twierdzenia
\ref{ma1:kh:thm:kh-primitive-derivative}.
Wobec zgodności sum KS i KH wniosek
\ref{ma1:kh:cor:ks-order-limits-monotone}
stosuje się zatem wprost do (a) oraz (b).
Daje całkowalność KH granicy, zbieżność całek
oraz $\int_{[a,b]}|f_n-f|\,\mathrm d\lambda_1\to0$.
W (c) stosujemy (a) do $-f_n\uparrow-f$, ponieważ
$\sup_n(\HK)\int_a^b(-f_n)=-\inf_n(\HK)\int_a^bf_n<\infty$.
Liniowość i niezmienność modułu przy zmianie znaku
dają te same wnioski dla $f_n,f$.

Z lokalności i liniowości KH oraz zgodności całek
Lebesgue'a i KH (twierdzenie \ref{it:thm:mcshane-lebesgue}
i inkluzja McShane'a w KH) dla każdego $x\in[a,b]$ mamy
\[
 F_n(x)-F(x)=\int_{[a,x]}(f_n-f)\,\mathrm d\lambda_1.
\]
Na mocy twierdzenia \ref{ma1:kh:thm:ac-lebesgue-primitive}
różnica ta jest bezwzględnie ciągła, a jej pochodna
jest równa $f_n-f$ prawie wszędzie.
Twierdzenie \ref{ma1:thm:ac-banach-separable} daje
\[
 V_a^b(F_n-F)=\int_{[a,b]}|f_n-f|\,\mathrm d\lambda_1.
\]
Ponieważ $F_n(a)-F(a)=0$, dla każdego $x$ mamy
$|F_n(x)-F(x)|\le V_a^b(F_n-F)$.
Z definicji normy Alexiewicza i liniowości całki
wynika $\|[f_n-f]\|_A=\|F_n-F\|_\infty$.
To daje wszystkie deklarowane równości i oszacowania.
\end{proof}
\begin{przyklad}[Konieczność założeń]
Brak wspólnej całkowalnej majoranty może prowadzić
do niezbieżności całek, jak pokazuje wcześniejszy
przykład $f_n=n\chi_{(0,1/n]}$.
Z kolei $g_n(x)=\min\{n,1/x\}$ dla $x>0$, $g_n(0)=0$, rośnie punktowo do funkcji równej $1/x$ na $(0,1]$, a $0$ w zerze, i, rozbijając całkę w $1/n$, otrzymujemy
\[
 \int_0^1 g_n=n\cdot\frac1n+
                 \int_{1/n}^1\frac{\mathrm dx}{x}
 =1+\log n\longrightarrow\infty.
\] Sama monotoniczność bez skończonej granicy całek nie wystarcza.
\end{przyklad}

\section{Porównanie zakresów całkowania}
\label{ma1:kh:sec:comparison}
\begin{twierdzenie}[Hierarchia całek względem długości]\label{ma1:kh:thm:integral-hierarchy}
Niech $a<b$. Wśród rzeczywistych funkcji o skończonych wartościach na $[a,b]$ (a nie klas równoważności modulo równość prawie wszędzie) zachodzą inkluzje
\[
 \mathcal R=\mathcal D_{\mathrm{Darboux}}
 \subsetneq\mathcal L^1=\mathcal M
 \subsetneq\HK=\mathcal P_{\mathrm{Perron}}=\mathcal D_*
 \subsetneq\mathcal D_{\mathrm{sz}}.
\]
Tutaj $\mathcal D_{\mathrm{Darboux}}$ oznacza klasę funkcji całkowalnych w sensie Darboux, $\mathcal D_*$ i $\mathcal D_{\mathrm{sz}}$ odpowiednio wąską i szeroką klasę Denjoya. Na wspólnych klasach funkcji całkowalnych wszystkie wymienione całki mają tę samą wartość.
\end{twierdzenie}
\begin{proof}
Pierwsza równość wynika z twierdzenia \ref{ma1:thm:riemann-darboux-equivalence}. Twierdzenie \atlasref{ME-08::thm:riemann-to-lebesgue}{ME-08: zgodność całek Riemanna i Lebesgue'a} daje inkluzję $\mathcal R\subset\mathcal L^1$ oraz zgodność wartości całek. Równość $\mathcal M=\mathcal L^1$ i zgodność wartości daje twierdzenie \ref{it:thm:mcshane-lebesgue}, zatem również $\mathcal R\subset\mathcal M$. Właściwość pierwszej inkluzji wykazuje funkcja Dirichleta (przykład \ref{ma1:kh:ex:dirichlet}): jest mierzalna, ograniczona i ma całkę Lebesgue'a zero, natomiast jest nieciągła w każdym punkcie, więc kryterium \ref{it:thm:rs-oscillation} (dla $\alpha(x)=x$) wyklucza całkowalność Riemanna. Inkluzja $\mathcal M\subset\HK$ wynika z faktu, że podział KH jest szczególnym wolnym podziałem; właściwość wykazuje pochodna z przykładu \ref{ma1:kh:ex:conditional-derivative}. Równość KH z wąską całką Denjoya wynika z wniosku \ref{it:cor:denjoy-hk}, a z całką Perrona z twierdzenia \ref{it:thm:perron-hk}. Jeżeli $f\in\mathcal D_*$, to jej pierwotna $F\in\AC G_*$ spełnia $F'=f$ prawie wszędzie. Ponieważ $\AC G_*\subset\AC G$ i zwykła pochodna jest w punktach istnienia również pochodną aproksymatywną, ta sama $F$ realizuje szeroką całkę Denjoya. Obie wartości są więc równe $F(b)-F(a)$, a jednoznaczność w szerszej klasie wynika z twierdzenia \ref{it:thm:acg-uniqueness}. Właściwość inkluzji wynika z przykładu \ref{ma1:kh:ex:denjoy-strict}. Wszystkie kontrprzykłady na $[0,1]$ przenosimy na dowolny $[a,b]$ przez rosnącą afiniczną bijekcję. Dla $T(t)=a+(b-a)t$ przenosimy funkcję $q$ z $[0,1]$ na $[a,b]$ jako $q\circ T^{-1}$. Całkowalność Riemanna i Lebesgue'a jest równoważna całkowalności funkcji przeniesionej przez afiniczną zmianę zmiennej i jej odwrotność, przy czym całka mnoży się przez $b-a$. Dla KH równoważność wynika z zastosowania do obu zmian zmiennej twierdzenia \ref{ma1:kh:tw:hk-podstawienie}, a dla szerokiej całki Denjoya z analogicznego zastosowania twierdzenia \ref{ma1:kh:thm:wide-denjoy-substitution}.
\end{proof}
\begin{wniosek}[Jednostajna granica w szerokim sensie Denjoya]\label{ma1:kh:cor:wide-denjoy-uniform}
Jeżeli funkcje $f_n$ są całkowalne w szerokim sensie Denjoya na $[a,b]$ i $f_n\to f$ jednostajnie, to $f$ jest całkowalna w tym sensie i całki zbiegają. Nie zakładamy ograniczoności samych funkcji; różnice $f-f_n$ są ograniczone dla dostatecznie dużych $n$. Ponadto, dla pierwotnych
$F_n(x)=(\mathrm D_{\mathrm{sz}})\int_a^x f_n$ oraz
$F(x)=(\mathrm D_{\mathrm{sz}})\int_a^x f$, dla wszystkich
dostatecznie dużych $n$ zachodzi
\[
 V_a^b(F-F_n)\le (b-a)\sup_{x\in[a,b]}|f(x)-f_n(x)|.
\]
W szczególności pierwotne zbiegają jednostajnie i
całki zbiegają jednostajnie względem obu końców podprzedziału.
\end{wniosek}
\begin{proof}
Wszystkie $f_n$ są mierzalne w sensie Lebesgue'a z dowodu twierdzenia \ref{it:thm:acg-uniqueness}, zatem $f$ również. Dla dostatecznie dużego $n$ funkcja $f-f_n$ jest ograniczona i mierzalna, więc należy do $\mathcal L^1$. Z twierdzenia \ref{ma1:kh:thm:ac-lebesgue-primitive} funkcja $f-f_n$ ma bezwzględnie ciągłą pierwotną, która należy do $\AC G$ i ma pochodną aproksymatywną równą $f-f_n$ prawie wszędzie. Definicja \ref{it:def:denjoy-broad} i liniowość szerokiej całki Denjoya dają całkowalność $f$. Dla tych indeksów mamy
\[
 \left|(\mathrm D_{\mathrm{sz}})\int_a^b(f-f_n)\right|
 \le (b-a)\|f-f_n\|_\infty\longrightarrow0,
\]
ponieważ różnica jest całkowalna w sensie Lebesgue'a.
Z liniowości i lokalności całki szerokiej oraz zgodności
z całką Lebesgue'a wynika dla każdego $x\in[a,b]$,
że $F(x)-F_n(x)=\int_a^x(f-f_n)\,\mathrm d\lambda_1$.
Dla $a<b$ twierdzenia \ref{ma1:kh:thm:ac-lebesgue-primitive}
i \ref{ma1:thm:ac-banach-separable} dają zatem
\[
 V_a^b(F-F_n)=\int_a^b|f-f_n|\,\mathrm d\lambda_1
 \le(b-a)\sup_{x\in[a,b]}|f(x)-f_n(x)|.
\]
Dla $a=b$ powyższe oszacowanie jest bezpośrednie.
Ponieważ $F(a)-F_n(a)=0$, moduł różnicy w każdym
punkcie jest ograniczony przez to wahanie.
Dla dowolnych $a\le u\le v\le b$ moduł całki
$(\mathrm D_{\mathrm{sz}})\int_u^v(f-f_n)$ jest
modułem przyrostu $F-F_n$ i również nie przekracza
tego samego wahania. Prawa strona oszacowania dąży
do zera niezależnie od $x,u,v$.
\end{proof}
\begin{uwaga}[Jak czytać hierarchię]
Twierdzenie \ref{ma1:kh:thm:integral-hierarchy} dotyczy całkowania względem długości. Nie umieszczamy w nim bez dodatkowych założeń całek Stieltjesa: dla ustalonego integratora $g$ zachodzi $\mathcal R(g)\subset\mathrm{KS}(g)$ ze zgodnością wartości (twierdzenie \ref{tw:ks-cauchy}), ale klasy zależą od $g$. Przy $g(t)=t$ są to odpowiednio Riemann i KH. Dla $g\in\BV$ całkowalność względem $\mu_{V_g}$ wystarcza do całkowalności KS (wniosek \ref{ma1:kh:cor:ks-variation-measure-l1}), lecz nie jest konieczna (przykład \ref{ma1:kh:ex:ks-nonsummable-spikes}); sama całkowalność względem $|\mu_g|$ może nie wystarczać (przykład \ref{ma1:kh:ex:ks-cancelled-spikes}). Przy zmianach wartości na zbiorze zerowym należy zatem zawsze określić miarę i konwencję całkowania. Zakres rozdziału obejmuje skalarne całki na przedziałach i opisane całki Stieltjesa, nie wszystkie uogólnienia całkowania (np. całki wielowymiarowe, wektorowe lub całkowanie na grupach).
\end{uwaga}

\paragraph{Literatura uzupełniająca.}
Dla całek Denjoya i Perrona: S. Saks, \emph{Theory of the Integral},
wyd.~2, 1937; R.~A. Gordon, \emph{The Integrals of Lebesgue,
Denjoy, Perron, and Henstock}, AMS, 1994.
Dla całki KH: A. Fonda, \emph{The Kurzweil--Henstock Integral for
Undergraduates}, Birkh\"auser, 2018.
Dla całki KS: M. Tvrd\'y, \emph{Kurzweil--Stieltjes Integration}, \url{https://users.math.cas.cz/~tvrdy/ks-dekk.pdf}.
Dla szerokiej całki Denjoya: Y. Kubota, \emph{A constructive definition
of the approximately continuous Denjoy integral}, Canad. Math. Bull.
15 (1972), 103--109.
Dowody podstawowych twierdzeń miarowych, w tym twierdzenia
Riesza--Markowa--Kakutaniego, istnienia miary Haara oraz
twierdzenia Daniella--Stone'a, należą do teorii miary i całki
(ME-21, ME-23 i ME-32); nie powtarzamy ich w tym tomie.
\endgroup
